\documentclass[12pt]{article}

% ------------------------------------------------
% Paquetes
% ------------------------------------------------
\usepackage[utf8]{inputenc}
\usepackage[T1]{fontenc}
\usepackage[spanish,es-nodecimaldot]{babel}
\usepackage{geometry}
\usepackage{setspace}
\usepackage{enumitem}
\usepackage{titlesec}
\usepackage{parskip}

% ------------------------------------------------
% Configuración
% ------------------------------------------------
\geometry{
    left=2.5cm,
    right=2.5cm,
    top=2.5cm,
    bottom=2.5cm
}

\onehalfspacing

\titleformat{\section}
{\large\bfseries}
{\thesection.}
{0.5em}
{}

\titleformat{\subsection}
{\normalsize\bfseries}
{\thesubsection.}
{0.5em}
{}

% ------------------------------------------------
% Documento
% ------------------------------------------------
\begin{document}

\begin{center}
    {\Large \textbf{Preguntas de investigación}}\\[0.5cm]
    {\large Investigación en Matemáticas Aplicadas y Computación}

    \begin{enumerate}
        \item Carlos Emiliano Segura Loera
        \item Valeria Hernandez Santos
        \item Nehiel Ramirez Hernandez
    
    \end{enumerate}
\end{center}

\vspace{0.5cm}

\section{Exploración del problema de investigación}

A partir de las preguntas exploratorias propuestas en el texto
\textit{Claves para plantear preguntas de investigación}, se analiza
el problema relacionado con el fenómeno de \textit{model collapse} y
la incorporación recurrente de contenido sintético en el entrenamiento
de modelos de lenguaje.

\subsection{¿En qué consiste exactamente el problema que estamos abordando?}

El problema consiste en estudiar qué ocurre cuando contenido generado
por modelos de lenguaje se incorpora de manera recurrente a los datos
utilizados para entrenar nuevos modelos. En particular, interesa
determinar si este proceso reduce la diversidad del contenido y aumenta
la repetición de patrones.

\subsection{¿En qué circunstancias surge el problema y a quiénes afecta?}

El problema puede surgir cuando contenido sintético se mezcla con datos
humanos y posteriormente se utiliza para entrenar nuevos modelos. Esto
es relevante para el desarrollo de modelos de lenguaje y para la calidad
y diversidad de la información disponible en Internet.

\subsection{¿Cómo se suscitó nuestro interés por el problema?}

El interés surge a partir del fenómeno denominado ``Internet muerto'' y
del aumento de contenido generado mediante inteligencia artificial. Esto
lleva a cuestionar qué puede ocurrir cuando dicho contenido comienza a
ser utilizado nuevamente como información de entrenamiento.

\subsection{¿Qué necesitamos averiguar?}

Necesitamos determinar cómo cambia el contenido generado cuando aumenta
la presencia de datos sintéticos en el entrenamiento, considerando tanto
la proporción de estos datos como el número de generaciones de
entrenamiento.

\subsection{¿Qué limitaciones tenemos para investigar el problema?}

La principal limitación es la disponibilidad de recursos
computacionales. Por ello, el estudio deberá centrarse en modelos y
conjuntos de datos que puedan ser utilizados dentro de los recursos
disponibles.

\subsection{¿Qué facetas del problema resultan más relevantes?}

Las principales dimensiones de interés son la diversidad léxica y
semántica, la repetición de patrones y la relación de estos fenómenos
con la proporción de datos sintéticos y el número de generaciones.

\subsection{¿Qué impacto o consecuencias esperamos que tenga la investigación?}

Se busca contribuir a una mejor comprensión de los efectos de reutilizar
contenido sintético en el entrenamiento de modelos de lenguaje,
particularmente sobre la diversidad y repetición del contenido generado.

\section{Preguntas de investigación}

A partir de la exploración anterior se establece como problema central
el posible efecto de la incorporación recurrente de datos sintéticos en
el entrenamiento de modelos de lenguaje.

\subsection{Pregunta principal}

\begin{quote}
\textbf{¿Cómo cambia la diversidad del contenido generado por modelos
de lenguaje cuando se incorporan repetidamente datos sintéticos al
proceso de entrenamiento?}
\end{quote}

Esta pregunta permite estudiar el fenómeno de \textit{model collapse}
desde una dimensión concreta y susceptible de análisis: los cambios en
la diversidad del contenido generado a través de sucesivas generaciones
de entrenamiento.

\subsection{Preguntas secundarias}

\begin{enumerate}[label=\textbf{\arabic*.}]
    \item \textbf{¿Cómo cambia la diversidad léxica y semántica del
    contenido generado conforme aumenta el número de generaciones de
    entrenamiento con datos sintéticos?}

    \item \textbf{¿Cómo influye la proporción de datos sintéticos
    utilizados durante el entrenamiento en la diversidad del contenido
    generado?}

    \item \textbf{¿Cómo cambia la frecuencia de patrones, frases o
    estructuras repetitivas conforme se incorporan datos sintéticos de
    manera recurrente?}

    \item \textbf{¿Qué relación existe entre la proporción de datos
    sintéticos y el número de generaciones de entrenamiento en los
    cambios observados en la diversidad del contenido?}

    \item \textbf{¿Qué características del contenido generado permiten
    identificar la aparición de un proceso asociado con el
    \textit{model collapse}?}
\end{enumerate}

\section{Relación con la hipótesis y el capitulado}

Las preguntas planteadas permiten delimitar los principales aspectos
del fenómeno que serán estudiados. La pregunta principal establece la
relación general entre el entrenamiento recurrente con datos sintéticos
y la diversidad del contenido generado.

Las preguntas secundarias permiten abordar específicamente la
diversidad léxica y semántica, la proporción de datos sintéticos, la
repetición de patrones y el efecto acumulativo de las generaciones de
entrenamiento.

A partir de estas preguntas será posible desarrollar el capitulado y
formular una hipótesis que pueda ser contrastada mediante un diseño
experimental y métricas de diversidad y repetición.

\end{document}